% Representación en R^3 del problema de maximización de la utilidad (rehace el
% Anexo 1 del tema original). Panel izquierdo: utilidad con un punto de saturación
% dentro del conjunto presupuestario (la restricción no se satura, lambda = 0).
% Panel derecho: utilidad monótona; el máximo condicionado está sobre la recta de
% balance (la restricción se satura, lambda > 0).
% Autor: Víctor Gutiérrez Marcos
\documentclass[tikz,border=4pt]{standalone}
% ==========================================================================
% tikz-tcee.tex — Estilos comunes de los gráficos TikZ del temario
% Autor: Víctor Gutiérrez Marcos
%
% Cada gráfico es un fichero standalone en fuentes/<tema>/graficos/:
%
%   \documentclass[tikz,border=4pt]{standalone}
%   % ==========================================================================
% tikz-tcee.tex — Estilos comunes de los gráficos TikZ del temario
% Autor: Víctor Gutiérrez Marcos
%
% Cada gráfico es un fichero standalone en fuentes/<tema>/graficos/:
%
%   \documentclass[tikz,border=4pt]{standalone}
%   % ==========================================================================
% tikz-tcee.tex — Estilos comunes de los gráficos TikZ del temario
% Autor: Víctor Gutiérrez Marcos
%
% Cada gráfico es un fichero standalone en fuentes/<tema>/graficos/:
%
%   \documentclass[tikz,border=4pt]{standalone}
%   \input{../../../latex/tikz-tcee}
%   \begin{document}
%   \begin{tikzpicture}[tcee]
%     \ejes{Q}{P}{6}{5}
%     \draw[curva1] (0.5,4.5) -- (5,0.8) node[etiqueta,right] {$D$};
%     \capa{2}{\draw[curva2] ...;}   % en el vídeo aparece en el paso 2
%   \end{tikzpicture}
%   \end{document}
%
% Paleta: la de la web (morado, dorado, salvia) más tres colores de apoyo
% distinguibles en blanco y negro por el trazo.
% ==========================================================================
\usepackage[T1]{fontenc}
\usepackage{newpxtext,newpxmath}
\usepackage{xcolor}
\usepackage{tikz,pgfplots}
\pgfplotsset{compat=1.18}
\usetikzlibrary{arrows.meta,calc,positioning,patterns,decorations.pathreplacing,intersections,shapes.misc,backgrounds}

\definecolor{tceemorado}{HTML}{5F2987}
\definecolor{tceedorado}{HTML}{B8860B}
\definecolor{tceeverde}{HTML}{3C7D3A}
\definecolor{tceeazul}{HTML}{1F5F99}
\definecolor{tceerojo}{HTML}{B22222}
\definecolor{tceegris}{HTML}{767171}
\definecolor{tceesalvia}{HTML}{E2EFD9}

% Capas para el vídeo: \capa{n}{…} dibuja su contenido solo a partir del paso n.
% En el PDF, la web y Word se ve todo; generar-video.py compila el gráfico una
% vez por paso con \def\tceepaso{k} para que cada elemento aparezca al narrarlo.
\providecommand{\tceepaso}{1000}
\newcommand{\capa}[2]{\ifnum#1>\tceepaso\relax\else#2\fi}

\tikzset{
  tcee/.style={line cap=round,line join=round,font=\small,>={Stealth[length=2.2mm]}},
  eje/.style={->,thick,black},
  curva1/.style={very thick,tceemorado},
  curva2/.style={very thick,tceeazul},
  curva3/.style={very thick,tceeverde},
  curva4/.style={very thick,tceedorado},
  curvanueva/.style={very thick,tceerojo},
  desplazada/.style={very thick,dashed},
  guia/.style={thin,dashed,tceegris},
  punto/.style={circle,fill=black,inner sep=1.3pt},
  etiqueta/.style={font=\small,inner sep=2pt},
  flecha/.style={->,thick,tceerojo},
  area/.style={fill=tceemorado,fill opacity=0.18,draw=none},
  area2/.style={fill=tceedorado,fill opacity=0.22,draw=none},
}

% \ejes{etiqueta x}{etiqueta y}{largo x}{alto y}
\newcommand{\ejes}[4]{%
  \draw[eje] (0,0) -- (#3,0) node[below left=2pt and 0pt] {#1};
  \draw[eje] (0,0) -- (0,#4) node[left] {#2};
  \node[below left] at (0,0) {$0$};}
% \proyeccion{(x,y)}{etq x}{etq y}: líneas guía a los ejes con etiquetas
\newcommand{\proyeccion}[3]{%
  \draw[guia] let \p1=#1 in (\x1,0) node[below] {#2} |- (0,\y1) node[left] {#3};}

%   \begin{document}
%   \begin{tikzpicture}[tcee]
%     \ejes{Q}{P}{6}{5}
%     \draw[curva1] (0.5,4.5) -- (5,0.8) node[etiqueta,right] {$D$};
%     \capa{2}{\draw[curva2] ...;}   % en el vídeo aparece en el paso 2
%   \end{tikzpicture}
%   \end{document}
%
% Paleta: la de la web (morado, dorado, salvia) más tres colores de apoyo
% distinguibles en blanco y negro por el trazo.
% ==========================================================================
\usepackage[T1]{fontenc}
\usepackage{newpxtext,newpxmath}
\usepackage{xcolor}
\usepackage{tikz,pgfplots}
\pgfplotsset{compat=1.18}
\usetikzlibrary{arrows.meta,calc,positioning,patterns,decorations.pathreplacing,intersections,shapes.misc,backgrounds}

\definecolor{tceemorado}{HTML}{5F2987}
\definecolor{tceedorado}{HTML}{B8860B}
\definecolor{tceeverde}{HTML}{3C7D3A}
\definecolor{tceeazul}{HTML}{1F5F99}
\definecolor{tceerojo}{HTML}{B22222}
\definecolor{tceegris}{HTML}{767171}
\definecolor{tceesalvia}{HTML}{E2EFD9}

% Capas para el vídeo: \capa{n}{…} dibuja su contenido solo a partir del paso n.
% En el PDF, la web y Word se ve todo; generar-video.py compila el gráfico una
% vez por paso con \def\tceepaso{k} para que cada elemento aparezca al narrarlo.
\providecommand{\tceepaso}{1000}
\newcommand{\capa}[2]{\ifnum#1>\tceepaso\relax\else#2\fi}

\tikzset{
  tcee/.style={line cap=round,line join=round,font=\small,>={Stealth[length=2.2mm]}},
  eje/.style={->,thick,black},
  curva1/.style={very thick,tceemorado},
  curva2/.style={very thick,tceeazul},
  curva3/.style={very thick,tceeverde},
  curva4/.style={very thick,tceedorado},
  curvanueva/.style={very thick,tceerojo},
  desplazada/.style={very thick,dashed},
  guia/.style={thin,dashed,tceegris},
  punto/.style={circle,fill=black,inner sep=1.3pt},
  etiqueta/.style={font=\small,inner sep=2pt},
  flecha/.style={->,thick,tceerojo},
  area/.style={fill=tceemorado,fill opacity=0.18,draw=none},
  area2/.style={fill=tceedorado,fill opacity=0.22,draw=none},
}

% \ejes{etiqueta x}{etiqueta y}{largo x}{alto y}
\newcommand{\ejes}[4]{%
  \draw[eje] (0,0) -- (#3,0) node[below left=2pt and 0pt] {#1};
  \draw[eje] (0,0) -- (0,#4) node[left] {#2};
  \node[below left] at (0,0) {$0$};}
% \proyeccion{(x,y)}{etq x}{etq y}: líneas guía a los ejes con etiquetas
\newcommand{\proyeccion}[3]{%
  \draw[guia] let \p1=#1 in (\x1,0) node[below] {#2} |- (0,\y1) node[left] {#3};}

%   \begin{document}
%   \begin{tikzpicture}[tcee]
%     \ejes{Q}{P}{6}{5}
%     \draw[curva1] (0.5,4.5) -- (5,0.8) node[etiqueta,right] {$D$};
%     \capa{2}{\draw[curva2] ...;}   % en el vídeo aparece en el paso 2
%   \end{tikzpicture}
%   \end{document}
%
% Paleta: la de la web (morado, dorado, salvia) más tres colores de apoyo
% distinguibles en blanco y negro por el trazo.
% ==========================================================================
\usepackage[T1]{fontenc}
\usepackage{newpxtext,newpxmath}
\usepackage{xcolor}
\usepackage{tikz,pgfplots}
\pgfplotsset{compat=1.18}
\usetikzlibrary{arrows.meta,calc,positioning,patterns,decorations.pathreplacing,intersections,shapes.misc,backgrounds}

\definecolor{tceemorado}{HTML}{5F2987}
\definecolor{tceedorado}{HTML}{B8860B}
\definecolor{tceeverde}{HTML}{3C7D3A}
\definecolor{tceeazul}{HTML}{1F5F99}
\definecolor{tceerojo}{HTML}{B22222}
\definecolor{tceegris}{HTML}{767171}
\definecolor{tceesalvia}{HTML}{E2EFD9}

% Capas para el vídeo: \capa{n}{…} dibuja su contenido solo a partir del paso n.
% En el PDF, la web y Word se ve todo; generar-video.py compila el gráfico una
% vez por paso con \def\tceepaso{k} para que cada elemento aparezca al narrarlo.
\providecommand{\tceepaso}{1000}
\newcommand{\capa}[2]{\ifnum#1>\tceepaso\relax\else#2\fi}

\tikzset{
  tcee/.style={line cap=round,line join=round,font=\small,>={Stealth[length=2.2mm]}},
  eje/.style={->,thick,black},
  curva1/.style={very thick,tceemorado},
  curva2/.style={very thick,tceeazul},
  curva3/.style={very thick,tceeverde},
  curva4/.style={very thick,tceedorado},
  curvanueva/.style={very thick,tceerojo},
  desplazada/.style={very thick,dashed},
  guia/.style={thin,dashed,tceegris},
  punto/.style={circle,fill=black,inner sep=1.3pt},
  etiqueta/.style={font=\small,inner sep=2pt},
  flecha/.style={->,thick,tceerojo},
  area/.style={fill=tceemorado,fill opacity=0.18,draw=none},
  area2/.style={fill=tceedorado,fill opacity=0.22,draw=none},
}

% \ejes{etiqueta x}{etiqueta y}{largo x}{alto y}
\newcommand{\ejes}[4]{%
  \draw[eje] (0,0) -- (#3,0) node[below left=2pt and 0pt] {#1};
  \draw[eje] (0,0) -- (0,#4) node[left] {#2};
  \node[below left] at (0,0) {$0$};}
% \proyeccion{(x,y)}{etq x}{etq y}: líneas guía a los ejes con etiquetas
\newcommand{\proyeccion}[3]{%
  \draw[guia] let \p1=#1 in (\x1,0) node[below] {#2} |- (0,\y1) node[left] {#3};}

\pgfplotsset{
  superficie/.style={surf,shader=faceted,fill opacity=0.45,draw=tceemorado!40,line width=0.15pt,
    colormap={tcee}{rgb255=(226,239,217) rgb255=(167,137,189)}},
  presupuesto/.style={fill=tceedorado,fill opacity=0.25,draw=tceedorado,thick},
  tresd/.style={width=7.6cm,height=7.2cm,view={125}{25},axis lines=none,
    xmin=0,xmax=4,ymin=0,ymax=4,zmin=0,zmax=3.6,clip=false},
}
\begin{document}
\begin{tikzpicture}[tcee]
  % ---------- panel izquierdo: restricción no operativa ----------
  \begin{axis}[tresd,name=a]
    % 1: ejes
    \capa{1}{\draw[eje] (axis cs:0,0,0) -- (axis cs:4.2,0,0) node[below] {$x_1$};}
    \capa{1}{\draw[eje] (axis cs:0,0,0) -- (axis cs:0,4.2,0) node[below] {$x_2$};}
    \capa{1}{\draw[eje] (axis cs:0,0,0) -- (axis cs:0,0,3.8) node[above] {$U(x)$};}
    % 2: conjunto presupuestario en el suelo
    \capa{2}{\addplot3[presupuesto] coordinates {(0,0,0) (3.5,0,0) (0,3.5,0) (0,0,0)};}
    % 3: utilidad con un máximo absoluto (punto de saturación)
    \capa{3}{\addplot3[superficie,domain=0:3.5,y domain=0:3.5,samples=12,samples y=12]
      {3*exp(-((x-1.2)^2+(y-1.2)^2)/2.5)};}
    % 4: el máximo, dentro del conjunto presupuestario
    \capa{4}{\draw[guia] (axis cs:1.2,1.2,3) -- (axis cs:1.2,1.2,0);}
    \capa{4}{\node[punto] at (axis cs:1.2,1.2,3) {};}
    \capa{4}{\node[punto] at (axis cs:1.2,1.2,0) {};}
    \capa{4}{\node[etiqueta,below] at (axis cs:1.2,1.2,0) {$x^*$};}
  \end{axis}
  \capa{4}{\node[font=\small,align=center,anchor=north] at ($(a.south)+(0,0.2)$)
    {(a) Máximo absoluto dentro\\ del conjunto presupuestario:\\ $p\cdot x^*<\overline{W}$, \ $\lambda=0$};}
  % ---------- panel derecho: restricción operativa ----------
  \begin{axis}[tresd,name=b,at={(a.east)},anchor=west,xshift=1.0cm]
    % 5: ejes
    \capa{5}{\draw[eje] (axis cs:0,0,0) -- (axis cs:4.2,0,0) node[below] {$x_1$};}
    \capa{5}{\draw[eje] (axis cs:0,0,0) -- (axis cs:0,4.2,0) node[below] {$x_2$};}
    \capa{5}{\draw[eje] (axis cs:0,0,0) -- (axis cs:0,0,3.8) node[above] {$U(x)$};}
    % 6: conjunto presupuestario
    \capa{6}{\addplot3[presupuesto] coordinates {(0,0,0) (3.5,0,0) (0,3.5,0) (0,0,0)};}
    % 7: utilidad monótona (siempre crece: no hay máximo absoluto)
    \capa{7}{\addplot3[superficie,domain=0:3.5,y domain=0:3.5,samples=12,samples y=12] {sqrt(x*y)};}
    % 8: utilidad a lo largo de la recta de balance (corte vertical)
    \capa{8}{\addplot3[curvanueva,domain=0:3.5,samples=40,samples y=0] (x,{3.5-x},{sqrt(x*(3.5-x))});}
    \capa{8}{\addplot3[guia,domain=0:3.5,samples=2,samples y=0] (x,{3.5-x},0);}
    % 9: máximo condicionado sobre la recta de balance
    \capa{9}{\draw[guia] (axis cs:1.75,1.75,1.75) -- (axis cs:1.75,1.75,0);}
    \capa{9}{\node[punto] at (axis cs:1.75,1.75,1.75) {};}
    \capa{9}{\node[punto] at (axis cs:1.75,1.75,0) {};}
    \capa{9}{\node[etiqueta,below] at (axis cs:1.75,1.75,0) {$x^*$};}
  \end{axis}
  \capa{9}{\node[font=\small,align=center,anchor=north] at ($(b.south)+(0,0.2)$)
    {(b) Utilidad monótona: máximo\\ condicionado sobre la recta de balance:\\ $p\cdot x^*=\overline{W}$, \ $\lambda>0$};}
\end{tikzpicture}
\end{document}
